\subsection{特異点が複数ある場合}

両端に特異性がある場合は,内部の1点$c$を固定して
\[
 \int_a^bf(x)\dd x
 =\int_a^cf(x)\dd x+\int_c^bf(x)\dd x
\]
と定める.ただし,右辺の2つの広義積分がそれぞれ収束することが必要である.両端を同じ速さで動かして相殺することは,この定義には含まれない.分割点$c$を別の内部の点へ変えても,その間の通常の定積分が一方に加わり,他方から引かれるだけなので,値は変わらない.

区間内部に特異点$d$があれば,\,$d$の左右にも分ける.また,実数全体の積分は,例えば
\[
 \int_{-\infty}^{\infty}f(x)\dd x
 =\int_{-\infty}^{0}f(x)\dd x+\int_0^{\infty}f(x)\dd x
\]
として,左右を独立に調べる.複数の特異点があるときは,それぞれに別の切り取り幅や上限を用意して考えると,何を確認すべきかが明確になる.

この違いを$\frac1x$で確かめよう.\,$\varepsilon>0$に対して
\[
 \int_{-1}^{-\varepsilon}\frac1x\dd x=\log\varepsilon\longrightarrow-\infty,
 \qquad
 \int_{\varepsilon}^{1}\frac1x\dd x=-\log\varepsilon\longrightarrow+\infty.
\]
したがって$\int_{-1}^1\frac1x\dd x$は広義積分として収束しない.しかし,同じ$\varepsilon$を使って足せば常に$0$になる.このような対称的な切り取り方による極限を,コーシーの主値という.
\[
 \operatorname{PV}\int_{-1}^1\frac1x\dd x
 =\lim_{\varepsilon\to0+}
 \left(\int_{-1}^{-\varepsilon}\frac1x\dd x
       +\int_{\varepsilon}^{1}\frac1x\dd x\right)=0.
\]
主値が存在しても,広義積分が収束するとは限らない.実際,左を$\varepsilon$,右を$2\varepsilon$で切れば,2つの積分の和は$\log\varepsilon-\log(2\varepsilon)=-\log2$になる.切り取り方を結び付けて得られた相殺だけでは,各側の収束を示したことにならない.

無限区間でも同じ注意が要る.例えば$\frac{x}{1+x^2}$は奇関数なので,任意の$R$に対して$\int_{-R}^R\frac{x}{1+x^2}\dd x=0$である.一方,右半分の積分は
\[
 \int_0^R\frac{x}{1+x^2}\dd x=\frac12\log(1+R^2)\longrightarrow\infty
\]
となり発散する.したがって$\int_{-\infty}^{\infty}\frac{x}{1+x^2}\dd x$は収束しない.奇関数であることを使う前に,広義積分の収束を確認する必要がある.



